% =====================================================================
% CHAPITRES DES SPRINTS - RAPPORT PFE
% Mise en place d'un pipeline automatisé d'extraction et de valorisation
% des rapports laboratoires : OCR, restitution décisionnelle avec un
% assistant conversationnel
% Auteur : Mohamed Aziz BENJEMAA - Poulina Group Holding
% =====================================================================

% Packages recommandés (à inclure dans le préambule du rapport principal) :
% \usepackage[utf8]{inputenc}
% \usepackage[T1]{fontenc}
% \usepackage[french]{babel}
% \usepackage{graphicx}
% \usepackage{array}
% \usepackage{longtable}
% \usepackage{booktabs}
% \usepackage{tabularx}
% \usepackage{multirow}
% \usepackage{hyperref}

% =====================================================================
% CHAPITRE 3 : SPRINT 1 - IMPLÉMENTATION DU FRONTEND
% =====================================================================
\chapter{Sprint 1 : Implémentation du frontend}

\section*{Introduction}
Ce premier sprint marque le démarrage de la phase de réalisation du projet. Il a pour
objectif principal de concevoir et de mettre en place l'ensemble des interfaces
graphiques de l'application avant le branchement des fonctionnalités métier. Cette
étape conditionne fortement l'expérience utilisateur (UX) du produit final, puisque
les écrans construits ici serviront de support à toutes les interactions ultérieures
(authentification, import de rapports, consultation des résultats, dashboard, etc.).
Nous présentons dans ce chapitre le backlog du sprint, les cas d'utilisation associés,
les diagrammes d'analyse et de conception, la réalisation des maquettes ainsi que la
rétrospective du sprint.

\section{Backlog du sprint 1}
Le backlog du sprint regroupe les fonctionnalités liées au développement de l'interface
utilisateur. Le tableau~\ref{tab:backlog-sprint1} présente les user stories couvertes
par ce sprint.

\begin{table}[h!]
\centering
\caption{Backlog du sprint 1}
\label{tab:backlog-sprint1}
\renewcommand{\arraystretch}{1.3}
\begin{tabularx}{\textwidth}{|c|l|X|}
\hline
\textbf{ID} & \textbf{Fonctionnalité} & \textbf{User Story} \\ \hline
US1 & Page de connexion / inscription &
En tant qu'utilisateur, je veux disposer d'écrans de connexion et d'inscription clairs
afin d'accéder à l'application en toute simplicité. \\ \hline
US2 & Page de profil &
En tant qu'utilisateur, je veux une interface de profil afin de consulter et modifier
mes informations personnelles. \\ \hline
US3 & Page d'importation de rapports &
En tant qu'utilisateur, je veux une interface intuitive de dépôt de fichiers PDF afin
d'importer mes rapports. \\ \hline
US4 & Page Historique &
En tant qu'utilisateur, je veux une page listant chronologiquement mes rapports afin de
les retrouver facilement. \\ \hline
US5 & Page Dashboard &
En tant qu'utilisateur / administrateur, je veux un tableau de bord visuel afin de
suivre l'activité globale du système. \\ \hline
US-NF1 & Thème clair / sombre &
En tant qu'utilisateur, je veux pouvoir basculer entre un mode clair et un mode sombre
pour adapter l'interface à mon confort visuel. \\ \hline
US-NF2 & Design responsive &
En tant qu'utilisateur, je veux que l'application reste utilisable sur tablette et
ordinateur, quel que soit l'écran. \\ \hline
\end{tabularx}
\end{table}

\section{Diagrammes de cas d'utilisation et descriptions détaillées}

\subsection{Diagramme de cas d'utilisation}
Le diagramme de cas d'utilisation du sprint 1 met en évidence les interactions
graphiques entre les acteurs (client et administrateur) et les différentes
interfaces de l'application.

\begin{figure}[h!]
\centering
% \includegraphics[width=0.85\textwidth]{figures/uc_sprint1.png}
\fbox{\parbox{0.6\textwidth}{\centering [Figure 3.1 : Diagramme de cas d'utilisation du sprint 1]}}
\caption{Diagramme de cas d'utilisation du sprint 1}
\label{fig:uc-sprint1}
\end{figure}

\subsection{Description textuelle détaillée}

\subsubsection{Cas d'utilisation : « S'authentifier »}
\begin{table}[h!]
\centering
\renewcommand{\arraystretch}{1.3}
\begin{tabularx}{\textwidth}{|l|X|}
\hline
\textbf{Titre} & S'authentifier \\ \hline
\textbf{Description} & Permet à un utilisateur d'accéder à son espace personnel via
le formulaire de connexion. \\ \hline
\textbf{Acteurs} & Client, Administrateur \\ \hline
\textbf{Pré-condition} & L'utilisateur dispose d'un compte valide. \\ \hline
\textbf{Post-condition} & L'utilisateur est redirigé vers son tableau de bord. \\ \hline
\textbf{Scénario nominal} &
1. L'utilisateur saisit son email et son mot de passe. \newline
2. Le système valide les identifiants. \newline
3. L'utilisateur est redirigé vers la page d'accueil correspondant à son rôle. \\ \hline
\textbf{Scénario alternatif} &
2.a. Identifiants invalides : un message d'erreur est affiché et la saisie est demandée
à nouveau. \\ \hline
\end{tabularx}
\caption{Description du cas d'utilisation « S'authentifier »}
\end{table}

\subsubsection{Cas d'utilisation : « Consulter le tableau de bord »}
\begin{table}[h!]
\centering
\renewcommand{\arraystretch}{1.3}
\begin{tabularx}{\textwidth}{|l|X|}
\hline
\textbf{Titre} & Consulter le tableau de bord \\ \hline
\textbf{Description} & L'utilisateur accède à la page d'accueil affichant les
indicateurs et raccourcis principaux. \\ \hline
\textbf{Acteurs} & Client, Administrateur \\ \hline
\textbf{Pré-condition} & L'utilisateur est authentifié. \\ \hline
\textbf{Post-condition} & Les widgets et indicateurs synthétiques sont chargés. \\ \hline
\textbf{Scénario nominal} &
1. L'utilisateur clique sur « Dashboard ». \newline
2. Le frontend appelle l'API et récupère les indicateurs. \newline
3. Les widgets s'affichent. \\ \hline
\textbf{Scénario alternatif} &
2.a. Erreur réseau : un message d'indisponibilité s'affiche et un bouton de rechargement
est proposé. \\ \hline
\end{tabularx}
\caption{Description du cas d'utilisation « Consulter le tableau de bord »}
\end{table}

\subsubsection{Cas d'utilisation : « Basculer le thème »}
\begin{table}[h!]
\centering
\renewcommand{\arraystretch}{1.3}
\begin{tabularx}{\textwidth}{|l|X|}
\hline
\textbf{Titre} & Basculer le thème (clair / sombre) \\ \hline
\textbf{Description} & L'utilisateur peut changer le thème graphique de l'application. \\ \hline
\textbf{Acteurs} & Client, Administrateur \\ \hline
\textbf{Pré-condition} & L'utilisateur est sur une page de l'application. \\ \hline
\textbf{Post-condition} & Le thème choisi est appliqué et persistant. \\ \hline
\textbf{Scénario nominal} &
1. L'utilisateur clique sur l'icône de thème. \newline
2. Le frontend applique la nouvelle classe CSS. \newline
3. Le choix est stocké côté client. \\ \hline
\textbf{Scénario alternatif} & --- \\ \hline
\end{tabularx}
\caption{Description du cas d'utilisation « Basculer le thème »}
\end{table}

\section{Analyse du sprint 1}
La phase d'analyse a pour but de représenter les échanges entre les composants
participants à chaque scénario. Pour le sprint 1, ces échanges concernent principalement
le navigateur, les vues Django (templates) et les fichiers statiques (CSS/JS/Bootstrap).

\begin{figure}[h!]
\centering
\fbox{\parbox{0.6\textwidth}{\centering [Figure 3.2 : Diagramme de communication --- Authentification]}}
\caption{Diagramme de communication : authentification}
\label{fig:comm-auth}
\end{figure}

\begin{figure}[h!]
\centering
\fbox{\parbox{0.6\textwidth}{\centering [Figure 3.3 : Diagramme de communication --- Dashboard]}}
\caption{Diagramme de communication : dashboard}
\label{fig:comm-dash}
\end{figure}

\section{Conception du sprint 1}
La conception s'appuie sur des diagrammes de séquence qui détaillent la chronologie
des appels entre l'utilisateur, le navigateur et les vues Django.

\subsection{Diagramme de séquence : authentification}
\begin{figure}[h!]
\centering
\fbox{\parbox{0.6\textwidth}{\centering [Figure 3.4 : Diagramme de séquence --- Authentification]}}
\caption{Diagramme de séquence : authentification}
\label{fig:seq-auth}
\end{figure}

\textbf{Description :} l'utilisateur saisit ses identifiants, le formulaire est envoyé
en POST à la vue \texttt{login\_view}, qui valide les données et redirige vers le
dashboard correspondant.

\subsection{Diagramme de séquence : chargement du dashboard}
\begin{figure}[h!]
\centering
\fbox{\parbox{0.6\textwidth}{\centering [Figure 3.5 : Diagramme de séquence --- Dashboard]}}
\caption{Diagramme de séquence : dashboard}
\label{fig:seq-dash}
\end{figure}

\textbf{Description :} le template récupère les indicateurs synthétiques (nombre de
rapports, statut, etc.) puis les affiche dans des cartes Bootstrap.

\section{Réalisation du sprint 1}
La réalisation a permis de mettre en place l'ensemble des écrans de l'application en
utilisant HTML5, CSS3, JavaScript et Bootstrap~5, avec un design moderne, responsive
et un support complet des modes clair et sombre.

\begin{figure}[h!]
\centering
\fbox{\parbox{0.6\textwidth}{\centering [Figure 3.6 : Page de connexion]}}
\caption{Interface : page de connexion}
\end{figure}

\begin{figure}[h!]
\centering
\fbox{\parbox{0.6\textwidth}{\centering [Figure 3.7 : Page d'inscription]}}
\caption{Interface : page d'inscription}
\end{figure}

\begin{figure}[h!]
\centering
\fbox{\parbox{0.6\textwidth}{\centering [Figure 3.8 : Page d'accueil (Dashboard)]}}
\caption{Interface : dashboard}
\end{figure}

\begin{figure}[h!]
\centering
\fbox{\parbox{0.6\textwidth}{\centering [Figure 3.9 : Page Historique]}}
\caption{Interface : historique des rapports}
\end{figure}

\begin{figure}[h!]
\centering
\fbox{\parbox{0.6\textwidth}{\centering [Figure 3.10 : Page Profil]}}
\caption{Interface : profil utilisateur}
\end{figure}

\begin{figure}[h!]
\centering
\fbox{\parbox{0.6\textwidth}{\centering [Figure 3.11 : Mode sombre]}}
\caption{Interface : mode sombre activé}
\end{figure}

\section{Rétrospective du sprint 1}
La rétrospective de ce premier sprint a permis d'identifier les points positifs et les
axes d'amélioration :
\begin{itemize}
\item \textbf{Points positifs :} maquettes livrées dans les délais, design cohérent
entre toutes les pages, support du mode sombre opérationnel, responsive validé sur
tablette et ordinateur.
\item \textbf{Points à améliorer :} certaines micro-animations devront être affinées
dans les sprints suivants ; la cohérence des composants (boutons, formulaires) sera
formalisée dans une bibliothèque interne.
\item \textbf{Décisions :} mise en place d'un guide de style minimal et factorisation
des blocs HTML répétitifs sous forme de \emph{template tags} Django.
\end{itemize}

\section*{Conclusion}
Ce premier sprint a permis de poser les fondations visuelles de l'application : toutes
les pages principales ont été conçues et intégrées sous Bootstrap~5, avec une
expérience cohérente sur les différents écrans. Le sprint suivant sera consacré à
l'implémentation du backend (Django) afin de connecter ces interfaces à la logique
métier et à la base de données.


% =====================================================================
% CHAPITRE 4 : SPRINT 2 - IMPLÉMENTATION DU BACKEND
% =====================================================================
\chapter{Sprint 2 : Implémentation du backend}

\section*{Introduction}
Après la mise en place des interfaces utilisateur dans le sprint précédent, ce
deuxième sprint est consacré à la construction de la partie serveur de l'application.
L'objectif est de mettre en place, à l'aide du framework Django, l'ensemble des
mécanismes liés à la gestion des comptes, à l'authentification, au contrôle d'accès
basé sur les rôles (RBAC), à la gestion du profil utilisateur ainsi qu'à la gestion
administrative des utilisateurs et des rapports.

\section{Backlog du sprint 2}
\begin{table}[h!]
\centering
\caption{Backlog du sprint 2}
\label{tab:backlog-sprint2}
\renewcommand{\arraystretch}{1.3}
\begin{tabularx}{\textwidth}{|c|l|X|}
\hline
\textbf{ID} & \textbf{Fonctionnalité} & \textbf{User Story} \\ \hline
US1 & Authentification &
En tant qu'utilisateur, je veux m'authentifier afin d'accéder à l'application en
toute sécurité. \\ \hline
US2 & Gestion du profil &
En tant qu'utilisateur, je veux gérer mon profil afin de consulter et modifier mes
informations personnelles. \\ \hline
US4 & Historique des rapports &
En tant qu'utilisateur, je veux consulter l'historique de mes rapports afin de
retrouver mes documents. \\ \hline
US5 & Recherche et filtrage &
En tant qu'utilisateur, je veux rechercher et filtrer les rapports afin d'accéder
rapidement à une information. \\ \hline
US-A1 & Gestion des utilisateurs (admin) &
En tant qu'administrateur, je veux activer, désactiver ou supprimer un compte
utilisateur afin de contrôler les accès. \\ \hline
US-A2 & Contrôle d'accès basé sur les rôles &
En tant que système, je veux distinguer les droits du client et de l'administrateur
afin de protéger les ressources. \\ \hline
\end{tabularx}
\end{table}

\section{Diagrammes de cas d'utilisation et descriptions détaillées}

\subsection{Diagramme de cas d'utilisation}
\begin{figure}[h!]
\centering
\fbox{\parbox{0.6\textwidth}{\centering [Figure 4.1 : Diagramme de cas d'utilisation du sprint 2]}}
\caption{Diagramme de cas d'utilisation du sprint 2}
\end{figure}

\subsection{Description textuelle détaillée}

\subsubsection{Cas d'utilisation : « Gérer son profil »}
\begin{table}[h!]
\centering
\renewcommand{\arraystretch}{1.3}
\begin{tabularx}{\textwidth}{|l|X|}
\hline
\textbf{Titre} & Gérer son profil \\ \hline
\textbf{Description} & L'utilisateur consulte et met à jour ses informations
personnelles. \\ \hline
\textbf{Acteurs} & Client, Administrateur \\ \hline
\textbf{Pré-condition} & L'utilisateur est authentifié. \\ \hline
\textbf{Post-condition} & Les informations sont enregistrées dans la base de données. \\ \hline
\textbf{Scénario nominal} &
1. L'utilisateur ouvre la page « Profil ». \newline
2. Il modifie les champs souhaités. \newline
3. Il valide la sauvegarde. \newline
4. Le backend met à jour la base et confirme l'opération. \\ \hline
\textbf{Scénario alternatif} &
3.a. Données invalides : la vue renvoie le formulaire avec les erreurs. \\ \hline
\end{tabularx}
\caption{Description du cas d'utilisation « Gérer son profil »}
\end{table}

\subsubsection{Cas d'utilisation : « Gérer les utilisateurs (admin) »}
\begin{table}[h!]
\centering
\renewcommand{\arraystretch}{1.3}
\begin{tabularx}{\textwidth}{|l|X|}
\hline
\textbf{Titre} & Gérer les utilisateurs \\ \hline
\textbf{Description} & L'administrateur active, désactive ou supprime un compte
utilisateur. \\ \hline
\textbf{Acteurs} & Administrateur \\ \hline
\textbf{Pré-condition} & L'administrateur est authentifié avec un rôle valide. \\ \hline
\textbf{Post-condition} & L'état du compte ciblé est mis à jour. \\ \hline
\textbf{Scénario nominal} &
1. L'administrateur ouvre la liste des utilisateurs. \newline
2. Il sélectionne un utilisateur. \newline
3. Il applique l'action souhaitée (activer / désactiver / supprimer). \newline
4. Le système confirme. \\ \hline
\textbf{Scénario alternatif} &
3.a. Échec de l'opération : un message d'erreur est affiché ; l'état du compte reste
inchangé. \\ \hline
\end{tabularx}
\caption{Description du cas d'utilisation « Gérer les utilisateurs »}
\end{table}

\subsubsection{Cas d'utilisation : « Rechercher et filtrer les rapports »}
\begin{table}[h!]
\centering
\renewcommand{\arraystretch}{1.3}
\begin{tabularx}{\textwidth}{|l|X|}
\hline
\textbf{Titre} & Rechercher et filtrer les rapports \\ \hline
\textbf{Description} & L'utilisateur recherche un rapport par mots-clés et applique
des filtres (statut, date). \\ \hline
\textbf{Acteurs} & Client, Administrateur \\ \hline
\textbf{Pré-condition} & L'utilisateur est authentifié et a accès à des rapports. \\ \hline
\textbf{Post-condition} & La liste filtrée est affichée. \\ \hline
\textbf{Scénario nominal} &
1. L'utilisateur saisit un mot-clé et/ou choisit des filtres. \newline
2. La vue applique les filtres sur le QuerySet. \newline
3. Les résultats paginés sont retournés. \\ \hline
\textbf{Scénario alternatif} &
3.a. Aucun résultat : un message « Aucun rapport trouvé » est affiché. \\ \hline
\end{tabularx}
\caption{Description du cas d'utilisation « Rechercher et filtrer les rapports »}
\end{table}

\section{Analyse du sprint 2}
\begin{figure}[h!]
\centering
\fbox{\parbox{0.6\textwidth}{\centering [Figure 4.2 : Diagramme de communication --- Authentification]}}
\caption{Diagramme de communication : authentification}
\end{figure}

\begin{figure}[h!]
\centering
\fbox{\parbox{0.6\textwidth}{\centering [Figure 4.3 : Diagramme de communication --- Gestion utilisateurs]}}
\caption{Diagramme de communication : gestion des utilisateurs}
\end{figure}

\section{Conception du sprint 2}

\subsection{Diagramme de séquence : authentification}
\begin{figure}[h!]
\centering
\fbox{\parbox{0.6\textwidth}{\centering [Figure 4.4 : Diagramme de séquence --- Authentification]}}
\caption{Diagramme de séquence : authentification}
\end{figure}

\textbf{Description :} le formulaire de connexion est envoyé à la vue
\texttt{login\_view} ; Django délègue l'authentification au backend, crée la session
puis redirige selon le rôle.

\subsection{Diagramme de séquence : modification du profil}
\begin{figure}[h!]
\centering
\fbox{\parbox{0.6\textwidth}{\centering [Figure 4.5 : Diagramme de séquence --- Profil]}}
\caption{Diagramme de séquence : modification du profil}
\end{figure}

\subsection{Diagramme de séquence : gestion des utilisateurs (admin)}
\begin{figure}[h!]
\centering
\fbox{\parbox{0.6\textwidth}{\centering [Figure 4.6 : Diagramme de séquence --- Gestion utilisateurs]}}
\caption{Diagramme de séquence : gestion des utilisateurs}
\end{figure}

\section{Réalisation du sprint 2}
Le backend a été développé en Django en respectant un découpage clair en applications
(\emph{accounts}, \emph{reports}, \emph{dashboard}, etc.). Les modèles \texttt{User}
et \texttt{Profile} ont été étendus pour gérer les rôles, et la sécurité a été
renforcée par les protections natives de Django (CSRF, XSS, injection SQL).

\begin{figure}[h!]
\centering
\fbox{\parbox{0.6\textwidth}{\centering [Figure 4.7 : Connexion fonctionnelle]}}
\caption{Interface : connexion connectée au backend}
\end{figure}

\begin{figure}[h!]
\centering
\fbox{\parbox{0.6\textwidth}{\centering [Figure 4.8 : Profil utilisateur (formulaire)]}}
\caption{Interface : profil utilisateur en mode édition}
\end{figure}

\begin{figure}[h!]
\centering
\fbox{\parbox{0.6\textwidth}{\centering [Figure 4.9 : Console admin --- gestion des comptes]}}
\caption{Interface : gestion des utilisateurs côté administrateur}
\end{figure}

\begin{figure}[h!]
\centering
\fbox{\parbox{0.6\textwidth}{\centering [Figure 4.10 : Recherche et filtres dans l'historique]}}
\caption{Interface : recherche et filtres dans l'historique}
\end{figure}

\section{Rétrospective du sprint 2}
\begin{itemize}
\item \textbf{Points positifs :} architecture Django propre et modulaire,
authentification et RBAC opérationnels, gestion fluide des sessions, intégration
réussie avec les templates du sprint 1.
\item \textbf{Points à améliorer :} la migration depuis SQLite vers SQL Server demande
une attention particulière (paramétrage des drivers ODBC) ; certains messages d'erreur
gagneraient à être harmonisés.
\item \textbf{Décisions :} formalisation d'un module \texttt{accounts} séparé pour les
profils, ajout d'un décorateur custom \texttt{@admin\_required} pour les vues
sensibles.
\end{itemize}

\section*{Conclusion}
Ce sprint a permis de doter l'application d'un socle backend robuste et sécurisé,
prêt à accueillir le pipeline OCR du sprint suivant. Les utilisateurs peuvent
désormais se connecter, gérer leur profil et accéder à un historique filtrable ;
l'administrateur dispose en parallèle d'outils pour superviser les comptes.


% =====================================================================
% CHAPITRE 5 : SPRINT 3 - IMPLÉMENTATION DE L'OCR
% =====================================================================
\chapter{Sprint 3 : Implémentation du pipeline OCR}

\section*{Introduction}
Le c\oe ur du projet réside dans l'extraction automatisée des données contenues dans
les rapports laboratoires. Ce sprint est consacré à la mise en place du pipeline OCR :
import des fichiers PDF, extraction de l'entête et des tableaux, validation des
données par l'administrateur, puis génération d'une synthèse PDF standardisée. Le
choix de la bibliothèque \texttt{pdfplumber} permet d'obtenir une fidélité supérieure
à un OCR classique sur les PDF natifs.

\section{Backlog du sprint 3}
\begin{table}[h!]
\centering
\caption{Backlog du sprint 3}
\label{tab:backlog-sprint3}
\renewcommand{\arraystretch}{1.3}
\begin{tabularx}{\textwidth}{|c|l|X|}
\hline
\textbf{ID} & \textbf{Fonctionnalité} & \textbf{User Story} \\ \hline
US3 & Importation de rapports PDF &
En tant qu'utilisateur, je veux importer un rapport PDF afin de l'exploiter dans
l'application. \\ \hline
US7 & Extraction OCR / pdfplumber &
En tant que système, je veux extraire automatiquement les champs du rapport via
\texttt{pdfplumber} afin d'éviter la saisie manuelle. \\ \hline
US6 & Validation administrative &
En tant qu'administrateur, je veux valider ou refuser un rapport extrait afin de
garantir la fiabilité des données. \\ \hline
US9 & Synthèse PDF générée &
En tant qu'utilisateur, je veux générer une synthèse PDF personnalisée afin
d'exporter les résultats de manière claire. \\ \hline
\end{tabularx}
\end{table}

\section{Diagrammes de cas d'utilisation et descriptions détaillées}

\subsection{Diagramme de cas d'utilisation}
\begin{figure}[h!]
\centering
\fbox{\parbox{0.6\textwidth}{\centering [Figure 5.1 : Diagramme de cas d'utilisation du sprint 3]}}
\caption{Diagramme de cas d'utilisation du sprint 3}
\end{figure}

\subsection{Description textuelle détaillée}

\subsubsection{Cas d'utilisation : « Importer un rapport PDF »}
\begin{table}[h!]
\centering
\renewcommand{\arraystretch}{1.3}
\begin{tabularx}{\textwidth}{|l|X|}
\hline
\textbf{Titre} & Importer un rapport PDF \\ \hline
\textbf{Description} & L'utilisateur dépose un rapport PDF qui sera traité par le
pipeline OCR. \\ \hline
\textbf{Acteurs} & Client \\ \hline
\textbf{Pré-condition} & L'utilisateur est authentifié ; le fichier est au format
PDF. \\ \hline
\textbf{Post-condition} & Le fichier est enregistré et placé en file d'attente
d'extraction. \\ \hline
\textbf{Scénario nominal} &
1. L'utilisateur sélectionne un fichier PDF. \newline
2. Il valide l'envoi. \newline
3. Le backend stocke le fichier et déclenche l'extraction. \\ \hline
\textbf{Scénario alternatif} &
2.a. Format invalide ou taille excessive : un message d'erreur est affiché et le
fichier est rejeté. \\ \hline
\end{tabularx}
\caption{Description du cas d'utilisation « Importer un rapport PDF »}
\end{table}

\subsubsection{Cas d'utilisation : « Extraire les données (OCR / pdfplumber) »}
\begin{table}[h!]
\centering
\renewcommand{\arraystretch}{1.3}
\begin{tabularx}{\textwidth}{|l|X|}
\hline
\textbf{Titre} & Extraire les données \\ \hline
\textbf{Description} & Le système extrait l'entête (client, code, dates...) et les
tableaux (essai, résultat, unité) du PDF. \\ \hline
\textbf{Acteurs} & Système (déclenchement manuel possible par l'administrateur) \\ \hline
\textbf{Pré-condition} & Un fichier PDF a été importé. \\ \hline
\textbf{Post-condition} & Les champs structurés sont stockés en base et associés au
rapport. \\ \hline
\textbf{Scénario nominal} &
1. Le pipeline ouvre le PDF avec \texttt{pdfplumber}. \newline
2. Il extrait l'entête. \newline
3. Il extrait les tableaux. \newline
4. Les données sont sauvegardées avec le statut « En attente de validation ». \\ \hline
\textbf{Scénario alternatif} &
1.a. PDF illisible ou scanné : le système marque le rapport en échec et notifie
l'utilisateur. \\ \hline
\end{tabularx}
\caption{Description du cas d'utilisation « Extraire les données »}
\end{table}

\subsubsection{Cas d'utilisation : « Valider ou refuser un rapport »}
\begin{table}[h!]
\centering
\renewcommand{\arraystretch}{1.3}
\begin{tabularx}{\textwidth}{|l|X|}
\hline
\textbf{Titre} & Valider ou refuser un rapport \\ \hline
\textbf{Description} & L'administrateur examine les données extraites et décide de
les accepter ou de les refuser. \\ \hline
\textbf{Acteurs} & Administrateur \\ \hline
\textbf{Pré-condition} & Un rapport est en attente de validation. \\ \hline
\textbf{Post-condition} & Le rapport est marqué « Validé » ou « Refusé ». \\ \hline
\textbf{Scénario nominal} &
1. L'administrateur ouvre le rapport. \newline
2. Il compare les données extraites avec le PDF source. \newline
3. Il clique sur « Valider » ou « Refuser ». \newline
4. Le système met à jour le statut et notifie l'utilisateur. \\ \hline
\textbf{Scénario alternatif} &
3.a. Données partiellement correctes : l'administrateur peut corriger manuellement
avant de valider. \\ \hline
\end{tabularx}
\caption{Description du cas d'utilisation « Valider ou refuser un rapport »}
\end{table}

\subsubsection{Cas d'utilisation : « Générer une synthèse PDF »}
\begin{table}[h!]
\centering
\renewcommand{\arraystretch}{1.3}
\begin{tabularx}{\textwidth}{|l|X|}
\hline
\textbf{Titre} & Générer une synthèse PDF \\ \hline
\textbf{Description} & Le système produit un PDF propre et standardisé regroupant les
informations clés d'un rapport. \\ \hline
\textbf{Acteurs} & Client, Administrateur \\ \hline
\textbf{Pré-condition} & Le rapport est validé. \\ \hline
\textbf{Post-condition} & Un fichier PDF est généré et téléchargeable. \\ \hline
\textbf{Scénario nominal} &
1. L'utilisateur clique sur « Exporter ». \newline
2. Le backend rend un template HTML/PDF. \newline
3. Le fichier est téléchargé. \\ \hline
\textbf{Scénario alternatif} &
2.a. Échec de génération : message d'erreur et nouvelle tentative possible. \\ \hline
\end{tabularx}
\caption{Description du cas d'utilisation « Générer une synthèse PDF »}
\end{table}

\section{Analyse du sprint 3}
\begin{figure}[h!]
\centering
\fbox{\parbox{0.6\textwidth}{\centering [Figure 5.2 : Diagramme de communication --- Import et extraction]}}
\caption{Diagramme de communication : import et extraction OCR}
\end{figure}

\begin{figure}[h!]
\centering
\fbox{\parbox{0.6\textwidth}{\centering [Figure 5.3 : Diagramme de communication --- Validation]}}
\caption{Diagramme de communication : workflow de validation}
\end{figure}

\section{Conception du sprint 3}

\subsection{Diagramme de séquence : import et extraction}
\begin{figure}[h!]
\centering
\fbox{\parbox{0.6\textwidth}{\centering [Figure 5.4 : Diagramme de séquence --- Import et extraction]}}
\caption{Diagramme de séquence : import et extraction}
\end{figure}

\textbf{Description :} l'utilisateur dépose un PDF ; la vue Django sauvegarde le
fichier, déclenche le service d'extraction, qui appelle \texttt{pdfplumber}, parse
l'entête et les tableaux, puis persiste les données en base.

\subsection{Diagramme de séquence : validation par l'administrateur}
\begin{figure}[h!]
\centering
\fbox{\parbox{0.6\textwidth}{\centering [Figure 5.5 : Diagramme de séquence --- Validation admin]}}
\caption{Diagramme de séquence : validation administrative}
\end{figure}

\subsection{Diagramme de séquence : génération de la synthèse PDF}
\begin{figure}[h!]
\centering
\fbox{\parbox{0.6\textwidth}{\centering [Figure 5.6 : Diagramme de séquence --- Synthèse PDF]}}
\caption{Diagramme de séquence : génération de la synthèse PDF}
\end{figure}

\section{Réalisation du sprint 3}
Le module d'extraction a été implémenté en Python. \texttt{pdfplumber} est utilisé
pour la lecture des PDF natifs, avec deux étapes principales : l'extraction de
l'entête (expressions régulières sur le texte brut) et l'extraction des tableaux
(détection de colonnes « Essai », « Résultat », « Unité »). Les résultats sont
stockés dans le modèle \texttt{Report} et ses lignes liées.

\begin{figure}[h!]
\centering
\fbox{\parbox{0.6\textwidth}{\centering [Figure 5.7 : Interface d'import PDF]}}
\caption{Interface : import d'un rapport PDF}
\end{figure}

\begin{figure}[h!]
\centering
\fbox{\parbox{0.6\textwidth}{\centering [Figure 5.8 : Résultats extraits affichés]}}
\caption{Interface : visualisation des données extraites}
\end{figure}

\begin{figure}[h!]
\centering
\fbox{\parbox{0.6\textwidth}{\centering [Figure 5.9 : File d'attente de validation côté admin]}}
\caption{Interface : workflow de validation par l'administrateur}
\end{figure}

\begin{figure}[h!]
\centering
\fbox{\parbox{0.6\textwidth}{\centering [Figure 5.10 : Synthèse PDF générée]}}
\caption{Interface : synthèse PDF générée}
\end{figure}

\section{Rétrospective du sprint 3}
\begin{itemize}
\item \textbf{Points positifs :} précision d'extraction très satisfaisante sur les
PDF natifs ; gain de temps notable par rapport à la saisie manuelle ; workflow de
validation clair pour l'administrateur.
\item \textbf{Points à améliorer :} certains PDF scannés ne sont pas reconnus par
\texttt{pdfplumber} ; un fallback Tesseract OCR pourra être envisagé dans une
itération future. La normalisation des unités (mg/L, g/L, etc.) est aussi à
améliorer.
\item \textbf{Décisions :} découpler le module d'extraction sous forme de service
Python autonome pour faciliter son évolution.
\end{itemize}

\section*{Conclusion}
Ce sprint constitue l'apport principal du projet : l'automatisation de la lecture
des rapports laboratoires est maintenant opérationnelle. L'utilisateur peut importer
un PDF, le système l'analyse automatiquement, l'administrateur valide les données et
une synthèse claire peut être téléchargée. Le sprint suivant exploitera ces données
au sein d'un entrepôt de données pour la restitution décisionnelle.


% =====================================================================
% CHAPITRE 6 : SPRINT 4 - DATA WAREHOUSE ET POWER BI
% =====================================================================
\chapter{Sprint 4 : Implémentation du Data Warehouse et Power BI}

\section*{Introduction}
La donnée extraite par le pipeline OCR doit désormais être valorisée à des fins
décisionnelles. Ce sprint a pour objectif de construire un entrepôt de données
(Data Warehouse) dans SQL Server, alimenté par un processus ETL réalisé avec SSIS,
puis de mettre en place un tableau de bord Power BI permettant à l'administrateur de
suivre les indicateurs clés du projet.

\section{Backlog du sprint 4}
\begin{table}[h!]
\centering
\caption{Backlog du sprint 4}
\label{tab:backlog-sprint4}
\renewcommand{\arraystretch}{1.3}
\begin{tabularx}{\textwidth}{|c|l|X|}
\hline
\textbf{ID} & \textbf{Fonctionnalité} & \textbf{User Story} \\ \hline
US10 & Processus ETL avec SSIS &
En tant que système, je veux mettre en place un processus ETL avec SSIS afin de
transformer et charger les données dans SQL Server. \\ \hline
US11 & Administration des données (SSMS) &
En tant qu'administrateur, je veux administrer et vérifier les données avec SSMS
afin de garantir leur cohérence. \\ \hline
US12 & Tableau de bord Power BI &
En tant qu'administrateur, je veux consulter un tableau de bord Power BI afin de
visualiser les indicateurs clés du projet. \\ \hline
US-NF3 & Filtrage interactif Power BI &
En tant qu'utilisateur, je veux filtrer le tableau de bord par date, statut,
utilisateur ou client afin d'analyser les données selon différents angles. \\ \hline
\end{tabularx}
\end{table}

\section{Diagrammes de cas d'utilisation et descriptions détaillées}

\subsection{Diagramme de cas d'utilisation}
\begin{figure}[h!]
\centering
\fbox{\parbox{0.6\textwidth}{\centering [Figure 6.1 : Diagramme de cas d'utilisation du sprint 4]}}
\caption{Diagramme de cas d'utilisation du sprint 4}
\end{figure}

\subsection{Description textuelle détaillée}

\subsubsection{Cas d'utilisation : « Exécuter le processus ETL »}
\begin{table}[h!]
\centering
\renewcommand{\arraystretch}{1.3}
\begin{tabularx}{\textwidth}{|l|X|}
\hline
\textbf{Titre} & Exécuter le processus ETL \\ \hline
\textbf{Description} & Le package SSIS extrait les données de la base opérationnelle,
les transforme et les charge dans le Data Warehouse. \\ \hline
\textbf{Acteurs} & Administrateur (ordonnanceur ou lancement manuel) \\ \hline
\textbf{Pré-condition} & La base opérationnelle est accessible ; les tables cibles du
DW existent. \\ \hline
\textbf{Post-condition} & Les tables de faits et de dimensions sont rafraîchies. \\ \hline
\textbf{Scénario nominal} &
1. L'administrateur lance le package SSIS. \newline
2. Les Data Flow Tasks extraient les données. \newline
3. Les transformations sont appliquées (nettoyage, dédoublonnage,
normalisation des unités). \newline
4. Les données sont chargées dans le DW. \\ \hline
\textbf{Scénario alternatif} &
3.a. Erreur de transformation : la tâche est mise en échec et un log SSISDB est
généré. \\ \hline
\end{tabularx}
\caption{Description du cas d'utilisation « Exécuter le processus ETL »}
\end{table}

\subsubsection{Cas d'utilisation : « Consulter le tableau de bord Power BI »}
\begin{table}[h!]
\centering
\renewcommand{\arraystretch}{1.3}
\begin{tabularx}{\textwidth}{|l|X|}
\hline
\textbf{Titre} & Consulter le tableau de bord Power BI \\ \hline
\textbf{Description} & L'administrateur explore les indicateurs et filtre les
données selon ses besoins. \\ \hline
\textbf{Acteurs} & Administrateur \\ \hline
\textbf{Pré-condition} & Le DW est alimenté ; Power BI est connecté à SQL Server. \\ \hline
\textbf{Post-condition} & Les visualisations reflètent les filtres appliqués. \\ \hline
\textbf{Scénario nominal} &
1. L'administrateur ouvre le rapport Power BI. \newline
2. Il applique des filtres (date, statut, client...). \newline
3. Les visuels se mettent à jour. \\ \hline
\textbf{Scénario alternatif} &
1.a. Connexion à la source impossible : un message d'erreur est affiché et un
rafraîchissement manuel est proposé. \\ \hline
\end{tabularx}
\caption{Description du cas d'utilisation « Consulter le tableau de bord Power BI »}
\end{table}

\section{Analyse du sprint 4}

\subsection{Architecture du Data Warehouse}
Le modèle adopté est un schéma en étoile (\emph{star schema}) avec une table de
faits centrale (\texttt{Fact\_Rapports}) entourée de dimensions :
\texttt{Dim\_Client}, \texttt{Dim\_Date}, \texttt{Dim\_Essai},
\texttt{Dim\_Utilisateur} et \texttt{Dim\_Statut}.

\begin{figure}[h!]
\centering
\fbox{\parbox{0.6\textwidth}{\centering [Figure 6.2 : Schéma en étoile du Data Warehouse]}}
\caption{Schéma en étoile du Data Warehouse}
\end{figure}

\subsection{Diagramme de communication : pipeline ETL}
\begin{figure}[h!]
\centering
\fbox{\parbox{0.6\textwidth}{\centering [Figure 6.3 : Diagramme de communication --- Pipeline ETL]}}
\caption{Diagramme de communication : pipeline ETL SSIS}
\end{figure}

\section{Conception du sprint 4}

\subsection{Diagramme de séquence : exécution du package SSIS}
\begin{figure}[h!]
\centering
\fbox{\parbox{0.6\textwidth}{\centering [Figure 6.4 : Diagramme de séquence --- ETL SSIS]}}
\caption{Diagramme de séquence : exécution du package SSIS}
\end{figure}

\textbf{Description :} le package SSIS exécute successivement les tâches
\emph{Source}, \emph{Data Conversion}, \emph{Lookup} sur les dimensions et
\emph{Destination} vers les tables du DW.

\subsection{Diagramme de séquence : consultation Power BI}
\begin{figure}[h!]
\centering
\fbox{\parbox{0.6\textwidth}{\centering [Figure 6.5 : Diagramme de séquence --- Consultation Power BI]}}
\caption{Diagramme de séquence : consultation du dashboard Power BI}
\end{figure}

\section{Réalisation du sprint 4}
La réalisation comprend la création du Data Warehouse, la conception du package SSIS
sous Visual Studio (SSDT), et la construction du rapport Power BI.

\begin{figure}[h!]
\centering
\fbox{\parbox{0.6\textwidth}{\centering [Figure 6.6 : Tables du Data Warehouse dans SSMS]}}
\caption{Capture SSMS : tables de faits et dimensions}
\end{figure}

\begin{figure}[h!]
\centering
\fbox{\parbox{0.6\textwidth}{\centering [Figure 6.7 : Package SSIS --- Control Flow]}}
\caption{Package SSIS : Control Flow}
\end{figure}

\begin{figure}[h!]
\centering
\fbox{\parbox{0.6\textwidth}{\centering [Figure 6.8 : Package SSIS --- Data Flow]}}
\caption{Package SSIS : Data Flow (extraction, transformation, chargement)}
\end{figure}

\begin{figure}[h!]
\centering
\fbox{\parbox{0.6\textwidth}{\centering [Figure 6.9 : Tableau de bord Power BI --- Vue globale]}}
\caption{Tableau de bord Power BI : vue globale}
\end{figure}

\begin{figure}[h!]
\centering
\fbox{\parbox{0.6\textwidth}{\centering [Figure 6.10 : Tableau de bord Power BI --- Filtres interactifs]}}
\caption{Tableau de bord Power BI : filtres interactifs}
\end{figure}

\section{Rétrospective du sprint 4}
\begin{itemize}
\item \textbf{Points positifs :} pipeline ETL stable, modélisation en étoile claire,
indicateurs Power BI pertinents (nombre de rapports par jour, taux de validation,
répartition par client...).
\item \textbf{Points à améliorer :} la fréquence de rafraîchissement du DW devra être
planifiée via SQL Server Agent ; certaines transformations gagneraient à être
factorisées dans des procédures stockées.
\item \textbf{Décisions :} mise en place d'un job d'ordonnancement quotidien et
documentation des KPI publiés dans Power BI.
\end{itemize}

\section*{Conclusion}
Ce sprint a permis de transformer la donnée brute extraite des PDF en une véritable
ressource décisionnelle, exploitable par les responsables du laboratoire. Le
sprint suivant ajoutera une couche d'interaction conversationnelle pour rendre
l'analyse encore plus accessible.


% =====================================================================
% CHAPITRE 7 : SPRINT 5 - IMPLÉMENTATION DU CHATBOT
% =====================================================================
\chapter{Sprint 5 : Implémentation du ChatBot}

\section*{Introduction}
Pour faciliter encore davantage l'accès à l'information, ce sprint a pour objectif
d'intégrer un assistant conversationnel (ChatBot) à l'application. L'utilisateur
pourra interroger ses rapports en langage naturel et obtenir des réponses
contextualisées, sans avoir à naviguer manuellement dans l'historique ou le
tableau de bord.

\section{Backlog du sprint 5}
\begin{table}[h!]
\centering
\caption{Backlog du sprint 5}
\label{tab:backlog-sprint5}
\renewcommand{\arraystretch}{1.3}
\begin{tabularx}{\textwidth}{|c|l|X|}
\hline
\textbf{ID} & \textbf{Fonctionnalité} & \textbf{User Story} \\ \hline
US8 & Chat interactif avec l'IA &
En tant qu'utilisateur, je veux interroger un assistant intelligent afin d'obtenir
rapidement des informations sur les rapports récents. \\ \hline
US-A3 & Analyse contextuelle &
En tant qu'utilisateur, je veux que l'IA prenne en compte mes rapports validés
afin de produire des réponses précises et personnalisées. \\ \hline
US-A4 & Intégration au dashboard &
En tant qu'utilisateur, je veux accéder au ChatBot directement depuis le dashboard
afin de fluidifier mon expérience. \\ \hline
\end{tabularx}
\end{table}

\section{Diagrammes de cas d'utilisation et descriptions détaillées}

\subsection{Diagramme de cas d'utilisation}
\begin{figure}[h!]
\centering
\fbox{\parbox{0.6\textwidth}{\centering [Figure 7.1 : Diagramme de cas d'utilisation du sprint 5]}}
\caption{Diagramme de cas d'utilisation du sprint 5}
\end{figure}

\subsection{Description textuelle détaillée}

\subsubsection{Cas d'utilisation : « Interroger l'assistant IA »}
\begin{table}[h!]
\centering
\renewcommand{\arraystretch}{1.3}
\begin{tabularx}{\textwidth}{|l|X|}
\hline
\textbf{Titre} & Interroger l'assistant IA \\ \hline
\textbf{Description} & L'utilisateur pose une question en langage naturel et reçoit
une réponse contextualisée à ses rapports. \\ \hline
\textbf{Acteurs} & Client, Administrateur \\ \hline
\textbf{Pré-condition} & L'utilisateur est authentifié et dispose d'au moins un
rapport validé. \\ \hline
\textbf{Post-condition} & L'assistant retourne une réponse formatée. \\ \hline
\textbf{Scénario nominal} &
1. L'utilisateur ouvre le widget ChatBot. \newline
2. Il pose une question (ex.: « Combien de rapports refusés ce mois-ci ? »). \newline
3. Le système prépare le contexte (données pertinentes). \newline
4. Le modèle d'IA répond. \newline
5. La réponse est affichée dans la conversation. \\ \hline
\textbf{Scénario alternatif} &
3.a. Aucun contexte pertinent : le bot répond qu'il n'a pas trouvé de données et
propose des reformulations. \\ \hline
\end{tabularx}
\caption{Description du cas d'utilisation « Interroger l'assistant IA »}
\end{table}

\subsubsection{Cas d'utilisation : « Analyser un rapport spécifique »}
\begin{table}[h!]
\centering
\renewcommand{\arraystretch}{1.3}
\begin{tabularx}{\textwidth}{|l|X|}
\hline
\textbf{Titre} & Analyser un rapport spécifique \\ \hline
\textbf{Description} & L'utilisateur demande à l'IA un résumé ou une explication
d'un rapport donné. \\ \hline
\textbf{Acteurs} & Client, Administrateur \\ \hline
\textbf{Pré-condition} & Le rapport est validé. \\ \hline
\textbf{Post-condition} & Une analyse synthétique est restituée. \\ \hline
\textbf{Scénario nominal} &
1. L'utilisateur sélectionne un rapport et clique sur « Analyser avec l'IA ». \newline
2. Le système injecte les données du rapport dans le contexte. \newline
3. L'IA produit une analyse. \\ \hline
\textbf{Scénario alternatif} &
2.a. Données incomplètes : l'analyse précise les champs manquants. \\ \hline
\end{tabularx}
\caption{Description du cas d'utilisation « Analyser un rapport spécifique »}
\end{table}

\section{Analyse du sprint 5}
\begin{figure}[h!]
\centering
\fbox{\parbox{0.6\textwidth}{\centering [Figure 7.2 : Diagramme de communication --- ChatBot]}}
\caption{Diagramme de communication : interaction avec le ChatBot}
\end{figure}

\section{Conception du sprint 5}

\subsection{Diagramme de séquence : envoi d'une question}
\begin{figure}[h!]
\centering
\fbox{\parbox{0.6\textwidth}{\centering [Figure 7.3 : Diagramme de séquence --- Question ChatBot]}}
\caption{Diagramme de séquence : envoi d'une question au ChatBot}
\end{figure}

\textbf{Description :} le widget envoie la question à la vue \texttt{chat\_api} ;
celle-ci construit le \emph{prompt} avec le contexte (rapports validés de
l'utilisateur), appelle le modèle d'IA, puis renvoie la réponse au front.

\subsection{Diagramme de séquence : analyse contextuelle d'un rapport}
\begin{figure}[h!]
\centering
\fbox{\parbox{0.6\textwidth}{\centering [Figure 7.4 : Diagramme de séquence --- Analyse contextuelle]}}
\caption{Diagramme de séquence : analyse contextuelle d'un rapport}
\end{figure}

\section{Réalisation du sprint 5}
Le ChatBot est intégré sous forme de widget flottant accessible depuis le dashboard
et la page des rapports. Côté serveur, une vue dédiée prépare un contexte
\emph{Retrieval-Augmented Generation} (RAG) léger à partir des rapports validés de
l'utilisateur, puis interroge le modèle d'IA. Les réponses sont historisées par
session.

\begin{figure}[h!]
\centering
\fbox{\parbox{0.6\textwidth}{\centering [Figure 7.5 : Widget ChatBot ouvert]}}
\caption{Interface : widget ChatBot ouvert}
\end{figure}

\begin{figure}[h!]
\centering
\fbox{\parbox{0.6\textwidth}{\centering [Figure 7.6 : Conversation avec l'IA]}}
\caption{Interface : conversation avec l'assistant}
\end{figure}

\begin{figure}[h!]
\centering
\fbox{\parbox{0.6\textwidth}{\centering [Figure 7.7 : Analyse contextuelle d'un rapport]}}
\caption{Interface : analyse contextuelle d'un rapport}
\end{figure}

\section{Rétrospective du sprint 5}
\begin{itemize}
\item \textbf{Points positifs :} qualité des réponses très satisfaisante sur les
questions de synthèse ; intégration fluide au dashboard ; expérience utilisateur
nettement améliorée.
\item \textbf{Points à améliorer :} certaines réponses peuvent encore « halluciner »
si le contexte est trop large ; un meilleur filtrage des rapports inclus dans le
\emph{prompt} doit être mis en place.
\item \textbf{Décisions :} limiter le contexte aux 10 derniers rapports validés et
ajouter un disclaimer indiquant que les réponses sont générées par une IA.
\end{itemize}

\section*{Conclusion}
Avec ce sprint, l'application franchit un palier important : l'utilisateur peut
désormais dialoguer naturellement avec ses propres données. Le dernier sprint sera
consacré à la phase de vérification, d'amélioration et de validation finale du
produit.


% =====================================================================
% CHAPITRE 8 : SPRINT 6 - VÉRIFICATION, AMÉLIORATION ET VALIDATION
% =====================================================================
\chapter{Sprint 6 : Vérification, amélioration et validation finale}

\section*{Introduction}
Ce dernier sprint a pour objectif de consolider l'ensemble du produit avant sa
livraison. Il regroupe les activités de tests (unitaires, fonctionnels, de bout en
bout), de revue de code, d'amélioration des performances et de l'expérience
utilisateur, ainsi que la recette finale auprès des parties prenantes du
laboratoire.

\section{Backlog du sprint 6}
\begin{table}[h!]
\centering
\caption{Backlog du sprint 6}
\label{tab:backlog-sprint6}
\renewcommand{\arraystretch}{1.3}
\begin{tabularx}{\textwidth}{|c|l|X|}
\hline
\textbf{ID} & \textbf{Fonctionnalité} & \textbf{User Story} \\ \hline
US-V1 & Tests fonctionnels &
En tant qu'équipe projet, nous voulons exécuter des scénarios de tests fonctionnels
afin de garantir le bon fonctionnement de toutes les user stories. \\ \hline
US-V2 & Tests de performance &
En tant qu'équipe projet, nous voulons mesurer les temps de réponse afin de garantir
les exigences non fonctionnelles. \\ \hline
US-V3 & Amélioration UX &
En tant qu'utilisateur final, je veux une interface affinée et des messages d'erreur
clairs afin de fluidifier mon usage. \\ \hline
US-V4 & Recette utilisateur &
En tant que partie prenante, je veux valider le produit avant la mise en production
afin de m'assurer qu'il répond à mes besoins. \\ \hline
\end{tabularx}
\end{table}

\section{Diagrammes de cas d'utilisation et descriptions détaillées}

\subsection{Diagramme de cas d'utilisation}
\begin{figure}[h!]
\centering
\fbox{\parbox{0.6\textwidth}{\centering [Figure 8.1 : Diagramme de cas d'utilisation du sprint 6]}}
\caption{Diagramme de cas d'utilisation du sprint 6}
\end{figure}

\subsection{Description textuelle détaillée}

\subsubsection{Cas d'utilisation : « Exécuter les tests »}
\begin{table}[h!]
\centering
\renewcommand{\arraystretch}{1.3}
\begin{tabularx}{\textwidth}{|l|X|}
\hline
\textbf{Titre} & Exécuter les tests \\ \hline
\textbf{Description} & L'équipe projet exécute les jeux de tests unitaires,
fonctionnels et de performance. \\ \hline
\textbf{Acteurs} & Équipe projet \\ \hline
\textbf{Pré-condition} & L'application est déployée en environnement de recette. \\ \hline
\textbf{Post-condition} & Les rapports de tests sont produits ; les anomalies sont
tracées. \\ \hline
\textbf{Scénario nominal} &
1. Lancement des tests automatisés. \newline
2. Exécution des tests manuels selon le plan. \newline
3. Consolidation des résultats. \\ \hline
\textbf{Scénario alternatif} &
3.a. Échec d'un test : création d'un ticket et correction dans le sprint courant. \\ \hline
\end{tabularx}
\caption{Description du cas d'utilisation « Exécuter les tests »}
\end{table}

\subsubsection{Cas d'utilisation : « Recetter le produit »}
\begin{table}[h!]
\centering
\renewcommand{\arraystretch}{1.3}
\begin{tabularx}{\textwidth}{|l|X|}
\hline
\textbf{Titre} & Recetter le produit \\ \hline
\textbf{Description} & Les parties prenantes valident le produit final. \\ \hline
\textbf{Acteurs} & Product Owner, Administrateur métier \\ \hline
\textbf{Pré-condition} & Tous les sprints précédents sont clos. \\ \hline
\textbf{Post-condition} & Le produit est accepté ou des actions correctives sont
identifiées. \\ \hline
\textbf{Scénario nominal} &
1. Démonstration des fonctionnalités. \newline
2. Validation des critères d'acceptation. \newline
3. Signature du PV de recette. \\ \hline
\textbf{Scénario alternatif} &
2.a. Non-conformité : ouverture d'un ticket correctif. \\ \hline
\end{tabularx}
\caption{Description du cas d'utilisation « Recetter le produit »}
\end{table}

\section{Analyse du sprint 6}

\subsection{Stratégie de tests}
La stratégie repose sur trois niveaux : tests unitaires (modules Python / Django),
tests d'intégration (pipeline OCR \(\rightarrow\) base \(\rightarrow\) ETL \(\rightarrow\)
Power BI) et tests fonctionnels (parcours utilisateur de bout en bout).

\begin{figure}[h!]
\centering
\fbox{\parbox{0.6\textwidth}{\centering [Figure 8.2 : Diagramme de communication --- Cycle de tests]}}
\caption{Diagramme de communication : cycle de tests}
\end{figure}

\section{Conception du sprint 6}

\subsection{Plan de tests}
Un plan de tests structuré a été élaboré, couvrant l'ensemble des user stories
prioritaires. Le tableau~\ref{tab:plan-tests} présente un extrait des cas testés.

\begin{table}[h!]
\centering
\caption{Extrait du plan de tests}
\label{tab:plan-tests}
\renewcommand{\arraystretch}{1.3}
\begin{tabularx}{\textwidth}{|c|X|X|c|}
\hline
\textbf{ID} & \textbf{Scénario} & \textbf{Résultat attendu} & \textbf{Statut} \\ \hline
T1 & Connexion avec identifiants valides & Redirection vers le dashboard & OK \\ \hline
T2 & Connexion avec mot de passe erroné & Message d'erreur explicite & OK \\ \hline
T3 & Import d'un PDF natif & Extraction automatique des champs & OK \\ \hline
T4 & Import d'un fichier non PDF & Rejet avec message d'erreur & OK \\ \hline
T5 & Validation d'un rapport par l'administrateur & Statut « Validé » et notification & OK \\ \hline
T6 & Génération de la synthèse PDF & PDF téléchargé conforme & OK \\ \hline
T7 & Rafraîchissement du DW via SSIS & Tables alimentées sans erreur & OK \\ \hline
T8 & Filtrage du dashboard Power BI & Visuels mis à jour & OK \\ \hline
T9 & Question simple au ChatBot & Réponse pertinente et contextuelle & OK \\ \hline
T10 & Bascule clair / sombre & Thème appliqué et persistant & OK \\ \hline
\end{tabularx}
\end{table}

\subsection{Diagramme de séquence : exécution d'un test fonctionnel}
\begin{figure}[h!]
\centering
\fbox{\parbox{0.6\textwidth}{\centering [Figure 8.3 : Diagramme de séquence --- Test fonctionnel]}}
\caption{Diagramme de séquence : exécution d'un test fonctionnel}
\end{figure}

\section{Réalisation du sprint 6}
La réalisation a porté sur l'exécution des tests, la correction des anomalies
mineures détectées, l'amélioration de quelques messages d'erreur, l'optimisation
des requêtes les plus lentes et la rédaction du guide utilisateur.

\begin{figure}[h!]
\centering
\fbox{\parbox{0.6\textwidth}{\centering [Figure 8.4 : Rapport de tests automatisés]}}
\caption{Rapport de tests automatisés}
\end{figure}

\begin{figure}[h!]
\centering
\fbox{\parbox{0.6\textwidth}{\centering [Figure 8.5 : Mesure des temps de réponse]}}
\caption{Mesure des temps de réponse de l'API}
\end{figure}

\begin{figure}[h!]
\centering
\fbox{\parbox{0.6\textwidth}{\centering [Figure 8.6 : Démonstration de recette]}}
\caption{Démonstration de recette auprès des utilisateurs}
\end{figure}

\section{Rétrospective du sprint 6}
\begin{itemize}
\item \textbf{Points positifs :} la majorité des cas de tests passent au premier
essai ; les utilisateurs métiers ont validé le produit lors de la recette ; le
produit est jugé conforme aux besoins exprimés.
\item \textbf{Points à améliorer :} la couverture des tests automatisés peut encore
être étendue ; certains tests d'intégration sont longs et gagneraient à être
parallélisés.
\item \textbf{Décisions :} produire un guide utilisateur final et planifier un suivi
post-mise en production.
\end{itemize}

\section*{Conclusion}
Ce sprint final clôture la phase de réalisation. Le produit a été testé,
amélioré puis recetté ; l'ensemble des objectifs initiaux est atteint. Le chapitre
suivant présentera la conclusion générale et les perspectives d'évolution du projet.
